% Standalone problem document, generated from the adjacent README.md.
% Compile with XeLaTeX. Mathematical content is maintained in Markdown.
\documentclass[11pt,a4paper]{article}
\usepackage[fontsize=10.5pt]{scrextend}
\usepackage[margin=22mm,headheight=15pt,headsep=9mm,footskip=11mm]{geometry}
\usepackage{fontspec}
\setmainfont{texgyrepagella}[Extension=.otf,UprightFont=*-regular,BoldFont=*-bold,ItalicFont=*-italic,BoldItalicFont=*-bolditalic]
\setsansfont{texgyreheros}[Extension=.otf,UprightFont=*-regular,BoldFont=*-bold,ItalicFont=*-italic,BoldItalicFont=*-bolditalic]
\setmonofont{texgyrecursor}[Extension=.otf,UprightFont=*-regular,BoldFont=*-bold,ItalicFont=*-italic,BoldItalicFont=*-bolditalic,Scale=MatchLowercase]
\usepackage{amsmath,amssymb}
\usepackage{unicode-math}
\setmathfont{texgyrepagella-math.otf}
\usepackage{xcolor,microtype,enumitem,needspace,fancyhdr}
\usepackage{bookmark}
\definecolor{ink}{HTML}{17344B}
\definecolor{accent}{HTML}{197D80}
\definecolor{muted}{HTML}{596773}
\hypersetup{colorlinks=true,linkcolor=accent,urlcolor=accent,pdftitle={MF-20 - Decidability
of zero hitting by a matrix
exponential},pdfauthor={Open Problems in Numerical Linear Algebra}}
\urlstyle{same}
\setlength{\parindent}{0pt}
\setlength{\parskip}{5pt plus 1pt minus 1pt}
\linespread{1.04}
\setlength{\emergencystretch}{3em}
\widowpenalty=10000
\clubpenalty=10000
\displaywidowpenalty=10000
\allowdisplaybreaks[1]
\setlist{nosep,leftmargin=1.5em}
\providecommand{\tightlist}{\setlength{\itemsep}{0pt}\setlength{\parskip}{0pt}}
\providecommand{\pandocbounded}[1]{#1}
\pagestyle{fancy}
\fancyhf{}
\fancyhead[L]{\sffamily\footnotesize\color{muted}OPEN PROBLEMS IN NUMERICAL LINEAR ALGEBRA}
\fancyhead[R]{\sffamily\bfseries\footnotesize\color{accent}MF-20}
\fancyfoot[L]{\parbox[b]{0.9\textwidth}{\sffamily\fontsize{8}{10}\selectfont\color{muted}Editorial ratings; a bounded search cannot certify that a problem remains unsolved.\\Literature check: 2026-09-10}}
\fancyfoot[R]{\sffamily\footnotesize\color{muted}\thepage}
\renewcommand{\headrulewidth}{0.3pt}
\renewcommand{\headrule}{\hbox to\headwidth{\color{accent}\leaders\hrule height \headrulewidth\hfill}}
\makeatletter
\renewcommand{\section}{\Needspace{5\baselineskip}\@startsection{section}{1}{0pt}{12pt plus 2pt minus 2pt}{4pt}{\sffamily\bfseries\large\color{ink}}}
\renewcommand{\subsection}{\Needspace{4\baselineskip}\@startsection{subsection}{2}{0pt}{9pt plus 1pt minus 1pt}{3pt}{\sffamily\bfseries\normalsize\color{ink}}}
\makeatother
\setcounter{secnumdepth}{0}
\begin{document}
{\sffamily\small\color{muted}Matrix functions and stability\par}
\vspace{3pt}
{\raggedright\hyphenpenalty=10000\exhyphenpenalty=10000\sffamily\bfseries\fontsize{21}{24}\selectfont\color{ink}Decidability
of zero hitting by a matrix exponential\par}
\vspace{7pt}
{\sffamily\small\textbf{Difficulty:} extreme\quad\textbf{Importance:} broadly
interesting\par}
{\sffamily\small\bfseries\color{accent}Status: Partially resolved\par}
\vspace{4pt}
{\color{accent}\hrule height 0.8pt}
\vspace{6pt}
\textbf{Rating rationale:} Exact zero detection with arbitrary algebraic
frequencies is an extreme decidability problem; certified
continuous-system verification gives broad significance.

\subsection{Problem statement}\label{problem-statement}

Let \(\overline{\mathbb Q}_{\mathbb R}\) denote the real algebraic
numbers. Is there a single algorithm which, given \(d\ge1\),
\(A\in\overline{\mathbb Q}_{\mathbb R}^{d\times d}\),
\(u,v\in\overline{\mathbb Q}_{\mathbb R}^d\), and an interval
\(I=[a,b]\) with rational \(0\le a\le b\) or \(I=[a,\infty)\) with
rational \(a\ge0\), always halts and decides whether

\[
u^T e^{tA}v=0\qquad\text{for some }t\in I?
\]

An algebraic real input is represented exactly by an integer polynomial
and a rational isolating interval. The exponential is
\(e^{tA}=\sum_{j\ge0}(tA)^j/j!\). The witness time is a real number, not
necessarily rational or algebraic. There is no diagonalizability
assumption.

This is the continuous Skolem problem in matrix form. The equivalent
scalar formulation asks whether the solution of a homogeneous linear
differential equation with constant algebraic coefficients and algebraic
initial data has a zero in \(I\). Bounded and unbounded intervals are
included in the same entry.

The problem is hyperplane reachability for a continuous linear dynamical
system. Exact tangential contact can defeat simple sign-change
detection, making the distinction between certified zero testing and
numerical sampling important for matrix-exponential computations.

\subsection{References}\label{references}

\begin{enumerate}
\def\labelenumi{\arabic{enumi}.}
\tightlist
\item
  V. Chonev, J. Ouaknine and J. Worrell, \emph{On the Skolem Problem for
  Continuous Linear Dynamical Systems}, ICALP 2016, article 100, §1
  (equivalent ODE and matrix formulations), §3, Theorem 7 (bounded case
  conditional on Schanuel's conjecture).
  \href{https://arxiv.org/pdf/1506.00695}{Full preprint, v3 (2016)}.
\item
  P. Bacik et al., \emph{A survey of the Skolem and Positivity Problems
  for linear recurrence sequences} (2026), §9.1, continuous-time
  analogues.
  \href{https://people.mpi-sws.org/~joel/publications/skolem_and_positivity_survey26.pdf}{Author-hosted
  PDF}.
\end{enumerate}

Status check (2026-09-10): the original full statement and current arXiv
record were checked, along with the 2026 survey. The known
Schanuel-conditional bounded result and partial frequency restrictions
do not give an unconditional algorithm for all inputs. Searches for
``continuous Skolem decidability 2025 2026'' and the exact source title
found no general resolution. Results on logical ``Skolem functions''
address a different problem. This is a bounded search report.

\subsection{Audit --- 2026-09-10}\label{audit-2026-09-10}

Independently rechecked
\href{https://arxiv.org/pdf/1506.00695}{Chonev--Ouaknine--Worrell, §1
and Theorem 7}. Added the exact unconditional subclass supporting
Partially resolved: \(d\le2\), \(I=[0,\infty)\), proved in P. C. Bell,
J.-C. Delvenne, R. M. Jungers and V. D. Blondel,
\href{https://perso.uclouvain.be/vincent.blondel/publications/10BDJ.pdf}{\emph{The
continuous Skolem-Pisot problem}}, \emph{Theoretical Computer Science}
411 (2010), 3625--3634, Theorem 10. The Schanuel-conditional result and
reductions to a bounded problem are separate evidence of progress. The
\href{https://people.mpi-sws.org/~joel/publications/skolem_and_positivity_survey26abs.html}{2026
survey's author abstract} and targeted later searches yielded no
unconditional general resolution. Both ratings are retained.
\end{document}
